% !TEX root = ../main.tex
\documentclass[../main.tex]{subfiles}
\begin{document}

\chapter{Estado del arte}
\label{chap:estado_del_arte}

Este capítulo repasa el estado del arte de las principales tecnologías involucradas en el proyecto: los sistemas de visión artificial embebida, las interfaces de alta velocidad empleadas para la transmisión de vídeo hacia un PC, los sensores de imagen CMOS y las metodologías de verificación de hardware digital.

\section{Sistemas de visión artificial embebida}

Los sistemas de visión artificial embebida han evolucionado notablemente en la última década. Las soluciones basadas en microcontroladores, permiten capturar y procesar imágenes de baja resolución con latencias aceptables. Sin embargo, cuando se requiere el procesamiento de flujos de vídeo en tiempo real a resoluciones \ac{VGA} o superiores, las \ac{FPGA}s se presentan como la alternativa más adecuada gracias a su capacidad de procesar datos en paralelo y al acceso directo a los pines de entrada/salida sin la mediación de un sistema operativo.

Entre los trabajos más representativos en este ámbito destaca el uso de plataformas Zynq de Xilinx, que combinan un procesador ARM con lógica programable y permiten implementar \textit{pipelines} de procesamiento de imagen con aceleración hardware. No obstante, para aplicaciones donde el objetivo principal es la adquisición y transmisión de imagen sin procesamiento local intensivo, una \ac{FPGA} pura sin parte de \ac{PS} como la Artix-7 ofrece menores costes y una complejidad de diseño más controlable.

\section{Interfaces de alta velocidad hacia PC}

La transmisión de flujos de imagen sin comprimir hacia un PC exige interfaces de alta velocidad. Las alternativas más comunes son USB~3.0, GigE Vision y PCIe. En el presente proyecto se utiliza, por requisito, la familia FT232H de FTDI (Módulo UM232H-B). Este chip ofrece un modo \textit{Synchronous FIFO} de 8 bits que permite alcanzar tasas de transferencia de hasta 40~MB/s sobre USB~2.0~\cite{ft232h_datasheet}, con una interfaz síncrona válida para su implementación en \ac{VHDL}.

El modo \textit{Synchronous FIFO} del FT232H opera con un reloj de 60~MHz generado por el propio chip. La FPGA actúa como maestro del bus: para transmitir datos (FPGA$\rightarrow$PC) activa una señal y presenta el dato en el bus de datos; para recibir datos (PC$\rightarrow$FPGA) activa otra señal diferente para que el FT232H controle el bus y a continuación aserta otra señal para consumir los bytes disponibles.
\section{Sensores de imagen CMOS}

El sensor MT9V111 de ON~Semiconductor es un sensor CMOS \ac{VGA} (640$\times$480) que entrega datos en formato YCbCr~4:2:2. Dispone de una \textbf{interfaz} de configuración \textbf{I2C} de dos hilos y una interfaz de \textbf{datos} paralela de 8~bits con señales de sincronismo de trama (\texttt{FVALID}), línea (\texttt{LVALID}) y reloj de píxel (\texttt{PIXCLK})~\cite{mt9v111_datasheet}. 

\section{Metodologías de verificación de hardware digital}


La verificación funcional ha evolucionado desde los bancos de pruebas \textit{ad hoc} hacia metodologías estructuradas que promueven sobre todo la reutilización y la automatización. En el ámbito de SystemVerilog, la referencia en la industria es \ac{UVM}~\cite{uvm}, que ofrece un entorno de verificación orientado a objetos con generación de estímulos aleatorios restringidos y cobertura funcional. Sin embargo, requiere un simulador con soporte completo de SystemVerilog, lo que la aleja de los flujos basados en VHDL.

Dentro del ecosistema VHDL han surgido varias alternativas. \ac{OSVVM}~\cite{osvvm} aporta \textit{constrained random}, cobertura funcional y modelos de \textit{scoreboard} sobre VHDL puro. cocotb~\cite{cocotb}, por su parte, permite escribir los bancos de pruebas en Python, separándolo del lenguaje de descripción de hardware. Frente a estas, \ac{UVVM}~\cite{uvvm} propone un enfoque estructurado y modular especialmente orientado a la verificación de interfaces: proporciona una biblioteca de utilidades (\texttt{uvvm\_util}) con funciones de \textit{logging}, gestión de alertas y comprobación de valores (checks), y un framework de \ac{VVC}s que abstrae cada agente de verificación y permite coordinarlos desde un secuenciador central.

Para este trabajo se ha optado por \ac{UVVM} por ser una metodología nativa de VHDL —el lenguaje en el que se describe todo el diseño—, lo que evita introducir SystemVerilog en el flujo y mantener así un único lenguaje en \ac{RTL} y el banco de pruebas. Su modelo de \ac{VVC}s encaja a priori de forma natural con la verificación de las interfaces del sistema, y su biblioteca de comprobaciones y alertas facilita la trazabilidad frente a los requisitos especificados.

\end{document}
